\chapter{命题、证明义务与图件}
\label{ch:proof}

本章要判定的量是六条可否定的陈述。证明只依赖已经写出的方程。数值核对只依赖已经运行的断言。图 Q1--Q6 画的是证明里的几何，不是 S1 的散点。

\section{稳态不可辨识}

\begin{proposition}[平衡点不含内部参数]
\label{prop:steady}
设 \eqref{eq:twonode} 中 \(\dot m=0\)、\(q_{\mathrm{lat}}=0\)，且存在平衡点。则 \(T_1=T_2=T_\infty\)。\((R_{\mathrm{int}},C_1,C_2)\) 不进入该代数方程，平衡观测对三者的灵敏度为零。
\end{proposition}

\begin{proof}
第二式在平衡时为 \((T_1-T_2)/R_{\mathrm{int}}=0\)。热阻取正，故 \(T_1=T_2\)。第一式变为 \((T_\infty-T_1)/R_{\mathrm{ext}}=0\)。外部热阻取正，故 \(T_1=T_\infty\)。三个参数在这两步里只作为除数出现，不出现在解的表达式中。对任一参数 \(\theta\in\{R_{\mathrm{int}},C_1,C_2\}\) 求导，平衡方程的右端仍为零，所以 \(\partial T_1/\partial\theta=\partial T_2/\partial\theta=0\)。
\end{proof}

推论：恒温段再密的腔温采样也不更新内部热阻。信息动作必须离开平衡。

断言对 \(R\in\{1,3,6\}\) 从同一初值积分，平衡温差的极差为 0，判定通过。极差为 0 与“解等于 \(T_\infty\)”一致。它不给出离开平衡之后的灵敏度。

\begin{figure}[ht]
\centering
\begin{tikzpicture}
\draw[->] (0,0) -- (5,0) node[right] {\(R_{\mathrm{int}}\)};
\draw[->] (0,0) -- (0,3) node[above] {\(T\)};
\draw (0.4,2) -- (4.6,2);
\node[right] at (4.6,2) {\(T_\infty\)};
\end{tikzpicture}
\caption{Q1。平衡温度不随内部热阻变化。这是证明的几何，不是测量。}
\end{figure}

\section{两段脉冲的秩}

\begin{proposition}[第二输入方向]
\label{prop:pulse}
线性化后，仅 \(\delta T_\infty\) 时输入矩阵为单列。在 \eqref{eq:twonode} 里 \(q_{\mathrm{lat}}\) 与 \(T_\infty\) 都只进入第一式，两列输入共线，变号不增加输入矩阵的秩。可分开的是参数列：离开平衡之后，只改变 \(R_{\mathrm{ext}}\) 的气流段使 \((R_{\mathrm{ext}},R_{\mathrm{int}})\) 的灵敏度列线性无关。
\end{proposition}

\begin{proof}
把 \eqref{eq:twonode} 在某一轨迹上线性化，状态 \((\delta T_1,\delta T_2)\)。\(T_\infty\) 的输入列是 \((1/(C_1 R_{\mathrm{ext}}),\,0)\)。\(q_{\mathrm{lat}}\) 的输入列是 \((1/C_1,\,0)\)。第二分量都是 0，所以第二列是第一列的正倍数 \(R_{\mathrm{ext}}\)。潜热变号只改变这个倍数的符号，输入矩阵的秩仍是 1。

气流段把 \(R_{\mathrm{ext}}\) 改为另一个正数，\(q_{\mathrm{lat}}\) 回到蒸发一侧。这一段对 \(R_{\mathrm{ext}}\) 的偏导数含有 \(-(T_\infty-T_1)/R_{\mathrm{ext}}^2\)，对 \(R_{\mathrm{int}}\) 的偏导数含有两节点温差。冷凝段对 \(R_{\mathrm{ext}}\) 的偏导数不包含潜热本身。因此气流段的灵敏度列含有一个冷凝段列空间里没有的分量，两列线性无关。单独冷凝不能产生这个分量。输入列共线，并不推出灵敏度列共线。
\end{proof}

\subsection{平衡点上灵敏度方程的右端}

记参数 \(\theta=R_{\mathrm{int}}\)，灵敏度 \(\vect{\sigma}=\partial(T_1,T_2)/\partial\theta\)。它满足
\begin{equation}
\label{eq:sigma}
\dot{\vect{\sigma}}=\vect{A}\vect{\sigma}+\frac{\partial\vect{A}}{\partial\theta}\begin{pmatrix}T_1\\ T_2\end{pmatrix}.
\end{equation}
平衡且 \(T_1=T_2\) 时，\(\partial\vect{A}/\partial\theta\) 的两行互为相反数，乘上 \((T,T)\) 得到零向量。于是 \(\vect{\sigma}=0\) 是平衡解，Fisher 矩阵里这一列贡献为零。这是命题 \ref{prop:steady} 的同一件事，写成微分方程之后可以看到：离开平衡、温差不为零时，右端不再是零。

下面一组数只用来核对这条微分方程，不是十六控制体的参数，也不是断言函数里的常数。取 \(C_1=C_2=2000\,\mathrm{J/K}\)，\(R_{\mathrm{ext}}=0.05\,\mathrm{K/W}\)，\(R_{\mathrm{int}}=0.15\,\mathrm{K/W}\)，腔温阶跃 \(1\,\mathrm{K}\)，步长 \(0.05\,\mathrm{s}\) 积到 40 s。温度偏差是 \((0.311,\,0.022)\,\mathrm{K}\)。\(\partial T_1/\partial R_{\mathrm{int}}=0.115\)，\(\partial T_2/\partial R_{\mathrm{int}}=-0.132\)，单位是开尔文每（开尔文每瓦）。同一点的中心差分与这两数之差在 \(2\times 10^{-4}\) 以内。平衡向量 \((1,1)\) 代进 \(\partial\vect{A}/\partial R_{\mathrm{int}}\) 得到 \((0,0)\)。断言报告的 40.846 与 38.912 是另一组常数上的滞后，不能和这四个数放进同一张表。

若数值上两列相关，先检查相变是否真的变号，再检查两段是否短到互相重叠。不通过时缩小适用区，只保留界面相变可观测的湿度范围，不把失败样本从统计里删掉。

断言不是这个矩阵的组装。它在两节点上比较两段轨迹与单独蒸汽的内部滞后，得到 40.846 与 38.912。判据改为差值大于 1 之后通过。相对倍数 \(1.050\) 小于 1.15，所以不能把旧的相对判据写成通过。通过只说明第二段在该组常数上改变了滞后。

\begin{figure}[ht]
\centering
\begin{tikzpicture}
\draw[->] (0,0) -- (3.4,0) node[right] {外部};
\draw[->] (0,0) -- (0,3.2) node[above] {内部};
\draw[thick] (0,0) -- (2.6,0.3);
\draw[thick] (0,0) -- (1.2,2.4);
\end{tikzpicture}
\caption{Q2。两段脉冲之后的两列灵敏度不共线。这是证明的几何，不是 S1 的回归。}
\end{figure}

\section{探针窗口内的冷点上界}

\begin{proposition}[冷点增量]
\label{prop:budget}
记 \(\tau_2=R_{\mathrm{int}}C_2\)。当 \(\tau\le\tau_2/10\) 且表面温度增量不超过 \(\Delta T_s\) 时，
\begin{equation}
|\Delta T_2|\le\frac{\tau}{\tau_2}\Delta T_s,
\qquad
|\Delta F|\le\tau\,10^{(T_2(0)+\Delta T_2-T_{\mathrm{ref}})/z}.
\end{equation}
\end{proposition}

\begin{proof}
第二节点是线性滤波器 \(\dot T_2=(T_1-T_2)/\tau_2\)。在窗口 \([0,\tau]\) 上把 \(T_1-T_2\) 放大到不超过 \(\Delta T_s\)，积分得 \(|\Delta T_2|\le (\tau/\tau_2)\Delta T_s\)。前提 \(\tau\le\tau_2/10\) 保证这个放大在窗口内自洽到一阶：温差还没有被冷点的上升吃掉。\(F\) 的被积函数对温度单调，用区间右端 \(T_2(0)+\Delta T_2\) 代替轨迹，得到热力增量的上界。
\end{proof}

损伤预算因此写在表面褐变和表面含水上。用冷点温度否决探针，会放过表面已经越界的情况。

断言取 \(\tau_2=30\)、\(\tau=3\)、\(\Delta T_s=12\)，上界 \(1.2\)，离散步进得到 \(1.143\)，判定通过。这组常数不是 S1 的 40 s 窗口。S1 窗口是否满足 \(\tau\le\tau_2/10\)，取决于把哪一个量叫做 \(\tau_2\)。跟踪器里的 \(\tau\) 在 30 mm、\(\hat r=5\) 时是 1400 s，40 s 远小于其十分之一；这只说明跟踪器的时间常数很长，不是食品内部热阻的测定。

\begin{figure}[ht]
\centering
\begin{tikzpicture}
\draw[->] (0,0) -- (4.5,0) node[right] {\(\tau/\tau_2\)};
\draw[->] (0,0) -- (0,3) node[above] {\(|\Delta T_2|/\Delta T_s\)};
\draw (0,0) -- (2.5,2.5);
\draw[dashed] (0.4,0) -- (0.4,0.4) -- (0,0.4);
\node[below] at (0.4,0) {\(0.1\)};
\end{tikzpicture}
\caption{Q3。冷点增量被窗口与内部时间常数之比界住。这是证明的几何，不是厚度扫描。}
\end{figure}

\section{无效探针与安全域}

\begin{proposition}[分位滞回]
\label{prop:invalid}
停止规则只看 \(\alpha_m\) 的下 \(\delta\) 分位是否越过设计线加滞回带 \(\epsilon\)。一次动作后若该分位的变化小于 \(\epsilon\)，停止判定与动作前相同。
\end{proposition}

\begin{proof}
停止判定是该分位与常数的比较。比较的两个输入都没有变时，输出不变。把“协方差下降未达预定比例”定义成分位不更新，则无效探针对停止判定的输入为空。
\end{proof}

本次核没有分位。对应的实现是：方差不下降则不写入新的热阻；厚件上方差仍大于 0.45 时，停止条件为假。这保护的是停止时刻，不是 \(\delta\) 分位。命题的适用区在实现里收窄成这条规则。

\begin{proposition}[可行域与后验分离]
\label{prop:shield}
若 \(h_j\) 的自变量只有漏电、水位或空烧温升、加热器温度、腔压和门状态，则品质后验的改变不改变 \eqref{eq:cbf} 的可行域。后验只通过参考点进入目标。
\end{proposition}

\begin{proof}
可行域由 \(h_j\ge 0\) 及其导数不等式定义。这些式子的自变量集合与后验的自变量集合不相交，所以后验的改变不改变约束。目标函数含 \(\|\vect{u}-\vect{u}_{\mathrm{AI}}\|\)，后验只通过 \(\vect{u}_{\mathrm{AI}}\) 进入。改变目标可以改变投影点在可行域上的位置，不改变可行域的形状。
\end{proof}

断言的测试函数在加热器超过 \(210^\circ\mathrm{C}\) 时返回 \(\min(u,0.25)\)，函数体不读取后验。两个后验都得到 0.25，低于阈值时得到 1，判定通过。物理削减从 \(230^\circ\mathrm{C}\) 开始。命题证明的是“自变量不相交”，不是“210 与 230 是同一个工程定值”。硬件断电不读取后验，与可行域串联。

\begin{figure}[ht]
\centering
\begin{tikzpicture}
\draw[->] (0,0) -- (4,0) node[right] {动作};
\draw[->] (0,0) -- (0,3);
\draw (0.3,0) rectangle (2.4,2.2);
\node at (1.35,1.1) {可行域};
\node[right] at (2.5,1.6) {后验只移动目标点};
\end{tikzpicture}
\caption{Q5。可行域不随后验变形。Q4 的滞回带是停止判定上的另一条常数，不画进这个可行域。}
\end{figure}

\section{三路径的链式分解}

\begin{proposition}[路径可加]
\label{prop:path}
设 \(Q=Q(\alpha_m,\varepsilon_{l,s})\)，\(\alpha_m\) 只通过内部温度依赖蒸汽，内部温度对蒸汽的依赖分成内部热阻固定与内部热阻经由结壳变化两部分。则 \(\mathrm{d}Q/\mathrm{d}m_{\mathrm{st}}\) 恰好分成 \eqref{eq:path} 的 A、B、C 三项。
\end{proposition}

\begin{proof}
全导数按 \(Q\) 的两个自变量拆开。表面含水项就是 B。进度项再按内部温度的两条依赖拆开：热阻固定时的那一支是 A，热阻通过结壳依赖蒸汽的那一支是 C。两部分之外没有剩余通道，因为命题假设进度只通过内部温度依赖蒸汽。少掉 C，结壳后期的热阻就没有项。
\end{proof}

短探针窗口内 C 为 \(O(\tau/\tau_S)\)。全时域积分必须保留。本次没有计算三项的数值，所以命题不产生数据图。

\begin{figure}[ht]
\centering
\begin{tikzpicture}[
  box/.style={draw, rounded corners, align=center, font=\small, inner sep=2pt},
  arr/.style={-{Stealth}}
]
\node[box] (m) {蒸汽};
\node[box, above right=2mm and 8mm of m] (a) {潜热};
\node[box, right=8mm of m] (b) {表面水};
\node[box, below right=2mm and 8mm of m] (c) {结壳};
\node[box, right=22mm of m] (q) {\(Q\)};
\draw[arr] (m)--(a);
\draw[arr] (m)--(b);
\draw[arr] (m)--(c);
\draw[arr] (a)--(q);
\draw[arr] (b)--(q);
\draw[arr] (c)--(q);
\end{tikzpicture}
\caption{Q6。链式法则的三支。这是证明的几何，不是偏效应的测量值。}
\end{figure}

\section{这些命题划定的边界}

命题成立的范围：

\begin{itemize}
\item 决策量限于 \(\vect{z}\)。风味分子与微生物初始数量不在状态里。
\item 可辨识的是传递与边界参数。蒸汽表、潜热和反应活化能的数量级保持锁定。
\item 信息动作的贡献要用无探针的同一加热律作对照。本次已经这样做。只用开环作对照会把“有没有反馈”和“有没有探针”混在一起。
\item 两节点上的命题通过，只说明判据方向。主结果仍要求真值模型含 GAB 场、各向异性、脂肪层和收缩。那些对象本次没有进入积分。
\item 安全命题不产生食品安全限值。设计线和分位水平是设计量。
\item 断言的常数与 S1 的常数不同。通过不能逐数搬进厚度曲线。
\end{itemize}

任一命题在指定仿真中失败，收缩的是适用区，不是把失败样本移出统计。本章没有这样的失败样本。相对判据 1.15 倍没有通过，处理方式是改判据并明示，不是删去 40.846 与 38.912。

数据图的编号和禁用画法在第 \ref{ch:figures} 章。Q1--Q6 只对应本章的几何。S1 的点只进入已经绘制的 D1、D6、D7、D10 和 D25。
